\section{仅隔离控制的解析基准}
\label{sec:s1}
固定 $c=c_0$ 时，维持 $I=\eta$ 唯一确定隔离控制函数，启动状态、控制时长和累计新增感染均可由易感者轨迹表示。退出后恢复常规控制，感染人数继续下降，因此退出时刻与清零时间仍须分别计算。

\subsection{启动、控制轨迹与退出}
在条件\eqref{eq:model:trigger}下，常规轨迹先后两次经过 $I=\eta$，分别位于感染人数增加和减少的阶段。本文使用第一次交点。
\begin{theorem}
\label{thm:s1:start}\label{thm:s1:trigger}
设条件\eqref{eq:model:trigger}成立，并令
\begin{equation}
z_\eta=-\frac{1}{S_c}\exp\!\left[-\frac{C_0-\eta}{\rho_1}\right]
=-\frac{S_0}{S_c}\exp\!\left[-\frac{S_0}{S_c}+\frac{\eta-I_0}{\rho_1}\right].
\label{eq:s1:lambert-argument}
\end{equation}
则 $z_\eta\in(-e^{-1},0)$，唯一的启动状态及相应时间为
\begin{equation}
S^*=-S_cW_{-1}(z_\eta),\qquad S_c<S^*<S_0,
\label{eq:s1:Sstar-lambert}
\end{equation}
\begin{equation}
t_1=\frac{1}{\beta_1^0}\int_{S^*}^{S_0}
\frac{dS}{S[I_0+\frac{\beta_2^0}{\beta_1^0}(S_0-S)+\rho_1\ln(S/S_0)]}>0.
\label{eq:s1:t1}
\end{equation}
这里 $W_{-1}$ 为 Lambert $W$ 函数的下实分支。
\end{theorem}
\begin{proof}
将 $I=\eta$ 代入式\eqref{eq:model:first-integral}，得
$(-S/S_c)e^{-S/S_c}=z_\eta$。又由峰值公式，
\[
z_\eta=-e^{-1}\exp\!\left[-\frac{I_{\max}^{no}-\eta}{\rho_1}\right]\in(-e^{-1},0).
\]
在此区间，$W_0(z)\in(-1,0)$、$W_{-1}(z)<-1$\cite{Corless1996}。前者给出的状态小于 $S_c$，对应感染人数下降后的交点。由于式\eqref{eq:model:routine-orbit}在 $S\in(S_c,S_0)$ 上关于 $S$ 严格递减，且端点感染人数分别为 $I_{\max}^{no}$ 与 $I_0$，该区间恰有一个交点，须取 $W_{-1}$。沿常规易感者方程分离变量即得式\eqref{eq:s1:t1}；积分区间上 $I(S)\ge I_0>0$，所以 $t_1$ 有限。
\end{proof}
若 $\eta=I_0$，可取 $t_1=0,S^*=S_0$。若 $\eta=I_{\max}^{no}$，两交点合并为 $S_c$，控制时长为零。若 $\eta<I_0$，初始状态已不满足感染上限。

启动时有 $S^*>S_c$，常规隔离比例不足以阻止感染人数增加。为保持 $I=\eta$，需立即提高隔离比例，并随易感者减少而调整。记
\begin{equation}
\bar S=\frac{(1-\beta)\gamma N}{\beta c_0},\qquad
\frac{\bar S}{S_c}=(1-\beta)(1-q_0)<1.
\label{eq:s1:barS}
\end{equation}
\begin{theorem}
\label{thm:s1:qc-structure}\label{thm:s1:platform}
在条件\eqref{eq:model:trigger}下，固定 $c(t)=c_0$。阈值控制期内保持 $I(t)=\eta$ 所需的隔离比例唯一为
\begin{equation}
q_c(S)=1-\frac{\gamma N}{\beta c_0S},\qquad S_c\le S\le S^*.
\label{eq:s1:q-control}
\end{equation}
相应的状态轨迹和控制时长为
\begin{equation}
S(t)=\bar S+(S^*-\bar S)e^{-c_0\eta(t-t_1)/N},\qquad I(t)=\eta,
\label{eq:s1:S-analytic}
\end{equation}
\begin{equation}
\Delta t=t_2-t_1=\frac{N}{c_0\eta}\ln\frac{S^*-\bar S}{S_c-\bar S}.
\label{eq:s1:t2}
\end{equation}
以式\eqref{eq:s1:S-analytic}代入式\eqref{eq:s1:q-control}定义 $q_c(t)$，则完整控制为
\begin{equation}
c(t)=c_0,\qquad q(t)=\begin{cases}
q_0,&t<t_1,\\
q_c(t),&t_1\le t\le t_2,\\
q_0,&t>t_2.
\end{cases}
\label{eq:s1:q-piecewise}
\end{equation}
它满足 $I(t)\le\eta$ 对所有 $t\ge0$ 成立。
\end{theorem}
\begin{proof}
由 $I=\eta>0$ 和 $\dot I=0$，式\eqref{eq:model:balance}唯一解出式\eqref{eq:s1:q-control}。代入 $S$ 方程得
\[
\dot S=-\frac{c_0\eta}{N}(S-\bar S),\qquad S(t_1)=S^*,
\]
故有式\eqref{eq:s1:S-analytic}。由 $S^*>S_c>\bar S$，轨迹严格下降并在式\eqref{eq:s1:t2}给出的有限时间首次到达 $S_c$。该区间上 $q_c(S_c)=q_0$、$q_0\le q_c(S)<1$，控制满足界限。

反过来，将如此得到的时间函数 $q_c(t)$ 代回原方程，式\eqref{eq:s1:S-analytic}与 $I=\eta$ 满足前两个方程及启动初值，其余仓室由线性方程或积分确定。由命题\ref{prop:model:positive}的解唯一性，该控制确实实现恒定感染，而不只是形式上满足 $\dot I=0$。启动前 $I<\eta$；退出后 $S<S_c$，故 $\dot I<0$，全程不超过阈值。
\end{proof}

\begin{corollaryn}
\label{cor:s1:qc-structure}
上述控制在 $t_1$ 从 $q_0$ 跳升至
\begin{equation}
q_{\max}=1-\frac{\gamma N}{\beta c_0S^*}>q_0,
\label{eq:s1:qmax}
\end{equation}
在 $(t_1,t_2)$ 严格递减，并在 $t_2$ 连续恢复为 $q_0$。若另有追踪隔离比例的上限 $q_0\le q_{\rm cap}<1$，则上述仅隔离控制可实施，当且仅当
\[
q_{\max}\le q_{\rm cap}
\quad\Longleftrightarrow\quad
\frac{\beta c_0(1-q_{\rm cap})S^*}{N}\le\gamma.
\]
该条件成立时，能力限制不改变上述控制轨迹、控制时长及成本公式。
\end{corollaryn}
\begin{proof}
由于 $q_c(S)$ 关于 $S$ 严格增加，$S^*>S_c$ 给出启动跳跃。沿控制轨迹，
\begin{equation}
\frac{dq_c(t)}{dt}=-\frac{\gamma\eta}{\beta}
\frac{S(t)-\bar S}{S(t)^2}<0.
\label{eq:s1:qprime}
\end{equation}
又 $q_c(S_c)=q_0$，故退出时连续。控制比例的最大值在启动时取得，因此全时段的能力限制等价于检查这一最大值；由式\eqref{eq:s1:qmax}整理即得上述等价式。当 $q_{\rm cap}=q_{\max}$ 时，控制只在启动时恰好达到能力上限，随后严格低于上限。
\end{proof}
若 $q_{\rm cap}<q_{\max}$，则在规定启动状态取 $q=q_{\rm cap}$ 仍有 $\dot I>0$，截断原控制函数会使感染人数在启动后超过阈值。其他参数及 $q_{\rm cap}$ 固定时，命题\ref{prop:s1:sensitivity}表明提高触发区间内的阈值会降低 $q_{\max}$，保持已有的可行性。降低阈值则可能使所需隔离比例超过能力上限。

将式\eqref{eq:s1:S-analytic}代入式\eqref{eq:s1:q-control}，得到以启动或退出时刻为参照的时间开环控制：
\begin{equation}
\begin{split}
q_c(t)
&=1-\frac{1}{\left[\dfrac{\beta c_0S^*}{\gamma N}+\beta-1\right]
 e^{-c_0\eta(t-t_1)/N}+1-\beta}\\
&=1-\frac{1}{\left[\dfrac{q_0}{1-q_0}+\beta\right]
 e^{c_0\eta(t_2-t)/N}+1-\beta}.
\end{split}
\label{eq:s1:q-time-explicit}
\end{equation}
两种时间表达沿解析轨迹与式\eqref{eq:s1:q-control}等价，均由名义模型预先确定。感染人数偏离阈值后的纠偏不包含在这一控制函数中。三阶段数值轨迹见图\ref{fig:scenario1:single-sim}，具体数值在第\ref{sec:numerics}节的仅隔离基准比较中统一展示。

\subsection{退出后的动态与累计新增感染}
退出后的常规轨迹给出剩余易感者的极限，并在 $\eta>1$ 时确定清零时间。累计新增感染由启动前、阈值控制期和退出后三段的感染流入相加得到。
\begin{theorem}
\label{thm:s1:tend}\label{thm:s1:post}
从 $(S_c,\eta)$ 恢复常规控制后，$S(t)$ 和 $I(t)$ 对 $t>t_2$ 均严格递减，且
\begin{equation}
I(t)\to0,\qquad
S(t)\to S_\infty=-S_cW_0\!\left(-e^{-1-\eta/\rho_1}\right)\in(0,S_c).
\label{eq:s1:Sinfty}
\end{equation}
若 $\eta>1$，则唯一的 $t_{\rm end}>t_2$ 满足 $I(t_{\rm end})=1$、$I'(t_{\rm end})<0$。记 $S_{\rm end}=S(t_{\rm end})$，有
\begin{equation}
S_{\rm end}=-S_cW_0\!\left(-e^{-1-(\eta-1)/\rho_1}\right),
\label{eq:s1:Send}
\end{equation}
\begin{equation}
t_{\rm end}=t_2+\frac{1}{\beta_1^0}\int_{S_{\rm end}}^{S_c}
\frac{dS}{S[\eta+\frac{\beta_2^0}{\beta_1^0}(S_c-S)+\rho_1\ln(S/S_c)]}.
\label{eq:s1:tend}
\end{equation}
\end{theorem}
证明见附录\ref{sec:proof:post-dynamics}。
若 $\eta=1$，则 $I(t_2)=1$ 但 $I'(t_2)=0$，不属于上述终点定义。若 $\eta<1$，退出后始终 $I<1$，也不存在上述终点。由于 $\delta_q>0$ 且隔离感染仓室的输入趋于零，另有 $I_q(t)\to0$。$S_q,R$ 单调有界，因此也存在极限。

当 $\eta>1$ 时，按第\ref{sec:cost}节的定义，将分别进入 $I$ 和 $I_q$ 的累计新增感染合计为 $\Itcum$，统计至 $t_{\rm end}$。
\begin{theorem}
\label{thm:s1:Itcum}
在前述仅隔离控制下，若 $\eta>1$，则
\begin{equation}
\begin{split}
\Itcum={}&\frac{\beta}{\beta+q_0(1-\beta)}(S_0-S^*+S_c-S_{\rm end})\\
&+\beta\left[S^*-S_c+\bar S\ln\frac{S^*-\bar S}{S_c-\bar S}\right].
\end{split}
\label{eq:s1:cumulative}
\end{equation}
将 $S_{\rm end}$ 替换为 $S_\infty$，即得到计算至无限时刻的累计新增感染。
\end{theorem}
\begin{proof}
前后常规阶段均有
\[
\frac{\beta c_0SI}{N}=\frac{\beta}{\beta+q_0(1-\beta)}(-\dot S),
\]
所以两段积分分别由 $S_0-S^*$、$S_c-S_{\rm end}$ 给出。阈值控制期内使用 $dt=-N\,dS/[c_0\eta(S-\bar S)]$，得到
\[
\int_{t_1}^{t_2}\frac{\beta c_0S\eta}{N}\,dt
=\beta\int_{S_c}^{S^*}\frac{S}{S-\bar S}\,dS
=\beta\left[S^*-S_c+\bar S\ln\frac{S^*-\bar S}{S_c-\bar S}\right].
\]
三段相加即得结论；无穷时刻的表达由极限得到。
\end{proof}
\subsection{阈值与常规接触率的影响}
\label{sec:scaling-theory}
阈值与常规接触率都会改变启动状态，后者还改变退出状态。控制时长随阈值严格下降，但对常规接触率的变化方向不固定；求导时须同时计入 $S^*$ 和 $S_c$ 对参数的依赖。

除特别说明外，以下偏导数均固定 $N,S_0,I_0$ 及其余独立参数，并在相应触发条件的内部讨论。$\eta$ 作为独立设计参数；若它随其他参数同步调整，相应总变化应按复合函数求导。在 $q_0=0$ 或 $\beta=1$ 的边界，涉及该参数的导数理解为允许参数侧的单侧导数。$S_c,S_q,q_c$ 中的下标均按变量定义使用，不表示求导。

\begin{lemma}
\label{lem:s1:Sstar-sensitivity}
在条件\eqref{eq:model:trigger}下，启动状态满足
\begin{equation}
\frac{\partial S^*}{\partial\eta}<0,\qquad
\frac{\partial S^*}{\partial c_0}>0,\qquad
\frac{\partial S^*}{\partial\beta}>0,\qquad
\frac{\partial S^*}{\partial q_0}<0.
\label{eq:s1:Sstar-signs}
\end{equation}
\end{lemma}
证明见附录\ref{sec:proof:start-sensitivity}。

提高阈值使控制开始时的易感者人数减小，由此产生干预负担与累计新增感染之间的权衡。
\begin{proposition}
\label{prop:threshold-tradeoff}\label{prop:s1:sensitivity}
对于仅隔离控制，在触发区间内有
\begin{equation}
\frac{\partial t_1}{\partial\eta}>0,\quad
\frac{\partial q_{\max}}{\partial\eta}<0,\quad
\frac{\partial\Delta t}{\partial\eta}<0,\quad
\frac{\partial J_q}{\partial\eta}<0,\quad
\frac{\partial J}{\partial\eta}<0.
\label{eq:threshold-tradeoff}
\end{equation}
若进一步有 $\eta>1$，则 $\partial\Itcum/\partial\eta>0$。计算至无限时刻的累计新增感染也随 $\eta$ 严格增加。
\end{proposition}
证明见附录\ref{sec:proof:threshold-tradeoff}。
提高阈值降低所需最大隔离比例、控制时长和成本，但在 $\eta>1$ 时增加计算至清零终点的累计新增感染。清零时间还包括退出后的下降过程，其变化需结合式\eqref{eq:s1:tend}判断。

常规接触率的比较须先确定触发范围。固定其余参数，当
\[
c_0\le\frac{\gamma N}{\beta(1-q_0)S_0}
\]
时，常规峰值为 $I_0$；超过该值后，式\eqref{eq:s1:Imax-no}给出
\begin{equation}
\frac{dI_{\max}^{no}}{dc_0}
=-\frac{\rho_1}{c_0}\ln\frac{S_c}{S_0}>0.
\label{eq:s1:Imax-no-c0-derivative}
\end{equation}
此时峰值从 $I_0$ 严格增加至 $I_0+\beta_2^0 S_0/\beta_1^0$。若
\[
I_0<\eta<I_0+\frac{\beta_2^0}{\beta_1^0}S_0,
\]
则唯一的临界常规接触率 $c_{0,\min}(\eta)$ 满足
\begin{equation}
I_{\max}^{no}\bigl(c_{0,\min}(\eta)\bigr)=\eta,
\qquad c_{0,\min}(\eta)>\frac{\gamma N}{\beta(1-q_0)S_0}.
\label{eq:s1:c0-min-def}
\end{equation}
正控制时长恰在 $c_0>c_{0,\min}(\eta)$ 时出现。

\begin{proposition}
\label{prop:contact-effects}\label{cor:s1:duration-sensitivity}
在触发条件内部，对于仅隔离控制，
\[
\frac{\partial t_1}{\partial c_0}<0,\qquad
\frac{\partial q_{\max}}{\partial c_0}>0,\qquad
\frac{\partial\Delta t}{\partial\beta}>0,\qquad
\frac{\partial\Delta t}{\partial q_0}<0.
\]
接触率对控制时长的影响由
\begin{equation}
\frac{\partial\Delta t}{\partial c_0}
=\frac{N}{c_0^2\eta}
\left\{\frac{S^*}{S^*-\bar S}
\left[1+\frac{S_c}{S^*-S_c}\ln\frac{S_0}{S^*}\right]
-\ln\frac{S^*-\bar S}{S_c-\bar S}\right\}
\label{eq:duration-c0}
\end{equation}
确定，其符号不能统一为正或负。固定 $I_0<\eta<I_0+\beta_2^0S_0/\beta_1^0$ 时，$\Delta t$ 在 $c_0\in(c_{0,\min},\infty)$ 内至少有一个极大值，但不据此断言极大值唯一。
\end{proposition}
证明见附录\ref{sec:proof:contact-effects}。
式\eqref{eq:duration-c0}的等价符号判据为
\begin{equation}
\frac{\partial\Delta t}{\partial c_0}\gtrless0
\quad\Longleftrightarrow\quad
\ln\frac{S_0}{S^*}\gtrless
\frac{S^*-S_c}{S_c}
\left[\frac{S^*-\bar S}{S^*}\ln\frac{S^*-\bar S}{S_c-\bar S}-1\right],
\label{eq:duration-c0-equivalent}
\end{equation}
等号亦相应成立。右端非正时，$\partial\Delta t/\partial c_0>0$。

较高接触率使干预更早启动并要求更大的隔离比例，控制时长则由起止状态与时间尺度共同决定，因而可以向两个方向变化。

人口规模的比较还取决于初值与阈值的缩放方式。记 $s=S/N,i=I/N,\theta=\eta/N$，以下固定归一化初值；固定绝对初始感染人数或固定绝对阈值时，比较条件不同。
\begin{proposition}
\label{prop:scaling}
对于仅隔离控制，固定 $s_0,i_0$ 和 $\beta,\gamma,c_0,q_0$。在阈值控制被触发时，$t_1,\Delta t,q_{\max},J$ 以及存在内部拐点时的 $\lambda$ 对 $N,\eta$ 的依赖仅通过 $\theta=\eta/N$ 表现。最终累计新增感染除以 $N$ 后亦不依赖 $N$。但是，使用绝对终点 $I=1$ 时，$t_{\rm end}$ 及计算至该终点的累计新增感染比例一般仍含 $N$。
\end{proposition}
\begin{proof}
常规阶段的归一化系统为
\[
\dot s=-c_0[\beta+q_0(1-\beta)]si,\qquad
\dot i=[\beta c_0(1-q_0)s-\gamma]i.
\]
它不含 $N$，触发条件为 $i=\theta$，所以 $t_1$ 与 $s^*=S^*/N$ 仅依赖上述固定量和 $\theta$。又
\[
s_c=\frac{\gamma}{\beta c_0(1-q_0)},\qquad
\bar s=\frac{(1-\beta)\gamma}{\beta c_0},\qquad
\Delta t=\frac1{c_0\theta}\ln\frac{s^*-\bar s}{s_c-\bar s}.
\]
例如，二次成本的归一化表达为
\begin{equation}
J=\frac{2}{c_0\theta}\int_{s_c}^{s^*}
\frac{(1-s_c/s)^2}{s-\bar s}\,ds,
\label{eq:dom:Jtheta}
\end{equation}
其中仅含归一化状态和固定参数。代入控制表达可得其余结论。无限时间的终点为 $i\to0$，同样不含 $N$；有限的操作性终点却为 $i=1/N$，因而保留绝对人口尺度。
\end{proof}

小初值与低阈值的联合极限，以及固定正初值下的有限极限见补充材料第\ref{sec:sup:baseline-scan}节，完整参数扫描也列于该节。
\subsection{隔离比例的时间形态}
\label{sec:s1:shape}
仅隔离控制的隔离比例严格下降，其下降速度在存在内部拐点时达到唯一最大值。以下在 $c=c_0$ 下确定这一时刻及其在控制期内的相对位置。

\begin{theorem}
\label{thm:s1:inflection}
在条件\eqref{eq:model:trigger}下，$q_c(t)$ 在 $(t_1,t_2)$ 内存在唯一拐点，当且仅当
\begin{equation}
S_c<2\bar S<S^*,
\label{eq:s1:inflection-condition}
\end{equation}
等价地，$\frac12<(1-\beta)(1-q_0)<-\frac12W_{-1}(z_\eta)$。
内部拐点存在时，
\begin{equation}
S(t_{\rm inf})=2\bar S,\quad
 t_{\rm inf}=t_1+\frac{N}{c_0\eta}\ln\frac{S^*-\bar S}{\bar S},\quad
\qinf=1-\frac1{2(1-\beta)}.
\label{eq:s1:inflection-values}
\end{equation}
此时 $q_0<\qinf<q_{\max}$，且隔离比例的下降速度在 $t_{\rm inf}$ 唯一达到最大值。
\end{theorem}
\begin{proof}
对式\eqref{eq:s1:qprime}求导，得
\begin{equation}
\frac{d^2q_c(t)}{dt^2}
=\frac{\gamma c_0\eta^2}{\beta N}
\frac{(S-\bar S)(2\bar S-S)}{S^3}.
\label{eq:s1:qsecond}
\end{equation}
由于 $S>\bar S$，其符号由 $2\bar S-S$ 决定。$S$ 从 $S^*$ 严格下降至 $S_c$，所以恰在式\eqref{eq:s1:inflection-condition}成立时穿过 $2\bar S$，且只穿过一次。将该状态代入式\eqref{eq:s1:S-analytic}及\eqref{eq:s1:q-control}，即得时间和隔离比例。条件\eqref{eq:s1:inflection-condition}蕴含 $\beta<1/2$，故所列 $\qinf$ 表达式有定义。二阶导数先负后正，而一阶导数始终为负，故 $|dq_c/dt|$ 在该时刻达到唯一最大值。
\end{proof}

内部拐点存在时，其隔离比例只依赖 $\beta$。拐点在控制期内的相对位置为
\begin{equation}
\lambda=\frac{t_{\rm inf}-t_1}{t_2-t_1}
=\frac{\ln[(S^*-\bar S)/\bar S]}
{\ln[(S^*-\bar S)/(S_c-\bar S)]}\in(0,1).
\label{eq:s1:lambda}
\end{equation}
其中拐点前、后的两段时长分别为
\begin{equation}
t_{\rm inf}-t_1=\frac{N}{c_0\eta}\ln\frac{S^*-\bar S}{\bar S},\qquad
 t_2-t_{\rm inf}=\frac{N}{c_0\eta}\ln\frac{\bar S}{S_c-\bar S}.
\label{eq:s1:seg-high}
\end{equation}
$\lambda$ 是拐点前时长占控制时长的比例，数值越大，下降最快的时刻在控制期内越靠后。绝对时刻 $t_{\rm inf}$ 还取决于启动时刻和控制时长。

若 $2\bar S\le S_c$，则 $q_c''(t)<0$ 对 $t\in(t_1,t_2)$ 成立；等号对应拐点到达控制终点。若 $2\bar S\ge S^*$，则 $q_c''(t)>0$ 对 $t\in(t_1,t_2)$ 成立；等号对应拐点到达控制起点。两种等号情形均无内部拐点。特别地，$\beta=1$ 时 $\bar S=0$，始终不存在内部拐点。

相对位置的辅助导数、正负导数构造及数值结果见补充材料第\ref{sec:sup:inflection}节。各参数敏感性的成立条件汇总于附录\ref{sec:appendix:proofs}的表\ref{tab:sensitivity}。
